\documentclass[10pt,a4paper]{amsart}
\usepackage{amsmath,amssymb,amsthm,graphicx}
% Retained minimal preamble for the Stacks Project chapter more-algebra.tex (More on Algebra).
% The chapter file pinned-original-source.tex carries only \input{preamble} on its first line, so no
% preamble text lives inside the chapter itself.  The Stacks Project uses one single shared preamble for
% every chapter; the snapshot stages it identically as sources/stacks-chow/preamble.tex and
% sources/stacks-etale-cohomology/preamble.tex.  The block below is copied mechanically, line for line,
% from sources/stacks-chow/preamble.tex lines 14-18 (the Xy-pic setup) and 154-217 (the native theorem
% declarations together with the mathematical macros).  Nothing was added, removed or redefined; the
% wrapper class and the amsmath/amssymb/amsthm/graphicx base packages are supplied by the assembler.
% Chapter: more-algebra.tex -- More on Algebra.
% For commutative diagrams we use Xy-pic
\usepackage[all]{xy}

% We use 2cell for 2-commutative diagrams.
\xyoption{2cell}

% Theorem environments.
%
\theoremstyle{plain}
\newtheorem{theorem}[subsection]{Theorem}
\newtheorem{proposition}[subsection]{Proposition}
\newtheorem{lemma}[subsection]{Lemma}

\theoremstyle{definition}
\newtheorem{definition}[subsection]{Definition}
\newtheorem{example}[subsection]{Example}
\newtheorem{exercise}[subsection]{Exercise}
\newtheorem{situation}[subsection]{Situation}

\theoremstyle{remark}
\newtheorem{remark}[subsection]{Remark}
\newtheorem{remarks}[subsection]{Remarks}

\numberwithin{equation}{subsection}

% Macros
%
\def\lim{\mathop{\mathrm{lim}}\nolimits}
\def\colim{\mathop{\mathrm{colim}}\nolimits}
\def\Spec{\mathop{\mathrm{Spec}}}
\def\Hom{\mathop{\mathrm{Hom}}\nolimits}
\def\Ext{\mathop{\mathrm{Ext}}\nolimits}
\def\SheafHom{\mathop{\mathcal{H}\!\mathit{om}}\nolimits}
\def\SheafExt{\mathop{\mathcal{E}\!\mathit{xt}}\nolimits}
\def\Sch{\mathit{Sch}}
\def\Mor{\mathop{\mathrm{Mor}}\nolimits}
\def\Ob{\mathop{\mathrm{Ob}}\nolimits}
\def\Sh{\mathop{\mathit{Sh}}\nolimits}
\def\NL{\mathop{N\!L}\nolimits}
\def\CH{\mathop{\mathrm{CH}}\nolimits}
\def\proetale{{pro\text{-}\acute{e}tale}}
\def\etale{{\acute{e}tale}}
\def\QCoh{\mathit{QCoh}}
\def\Ker{\mathop{\mathrm{Ker}}}
\def\Im{\mathop{\mathrm{Im}}}
\def\Coker{\mathop{\mathrm{Coker}}}
\def\Coim{\mathop{\mathrm{Coim}}}

% Boxtimes
%
\DeclareMathSymbol{\boxtimes}{\mathbin}{AMSa}{"02}

%
% Macros for moduli stacks/spaces
%
\def\QCohstack{\mathcal{QC}\!\mathit{oh}}
\def\Cohstack{\mathcal{C}\!\mathit{oh}}
\def\Spacesstack{\mathcal{S}\!\mathit{paces}}
\def\Quotfunctor{\mathrm{Quot}}
\def\Hilbfunctor{\mathrm{Hilb}}
\def\Curvesstack{\mathcal{C}\!\mathit{urves}}
\def\Polarizedstack{\mathcal{P}\!\mathit{olarized}}
\def\Complexesstack{\mathcal{C}\!\mathit{omplexes}}
% \Pic is the operator that assigns to X its picard group, usage \Pic(X)
% \Picardstack_{X/B} denotes the Picard stack of X over B
% \Picardfunctor_{X/B} denotes the Picard functor of X over B
\def\Pic{\mathop{\mathrm{Pic}}\nolimits}
\def\Picardstack{\mathcal{P}\!\mathit{ic}}
\def\Picardfunctor{\mathrm{Pic}}
\def\Deformationcategory{\mathcal{D}\!\mathit{ef}}

% --- appended 2026-10-12 (Chief): the project own definitions of the reference/slogan/history
% display environments, copied mechanically from the shared Stacks preamble lines 8-12
% (\usepackage{verbatim} plus the three \newenvironment lines, which the chapter preamble keeps
% verbatim from the project preamble). Nothing added, removed or redefined: these environments
% carry no mathematical content and the project itself typesets them as comments. Without these
% three lines a counted statement whose own text uses them fails with Environment slogan undefined.
\usepackage{verbatim}
\newenvironment{reference}{\comment}{\endcomment}
\newenvironment{slogan}{\comment}{\endcomment}
\newenvironment{history}{\comment}{\endcomment}
\usepackage{cleveref}
\crefname{originallemma}{original lemma}{original lemmas}
\Crefname{originallemma}{Original lemma}{Original lemmas}
\makeatletter
\newlabel{lemma-derived-complete-bounded}{{0H83}{source}{}{sourcepointer-1}{}}
\newlabel{lemma-derived-complete-bounded@cref}{{[originallemma][0][]0H83}{[1][0][]source}{}{}{}}
\makeatother
\crefname{originallemma}{original lemma}{original lemmas}
\Crefname{originallemma}{Original lemma}{Original lemmas}
\makeatletter
\newlabel{lemma-derived-vanishing-mod-I}{{0G1T}{source}{}{sourcepointer-2}{}}
\newlabel{lemma-derived-vanishing-mod-I@cref}{{[originallemma][0][]0G1T}{[1][0][]source}{}{}{}}
\makeatother
\crefname{originallemma}{original lemma}{original lemmas}
\Crefname{originallemma}{Original lemma}{Original lemmas}
\makeatletter
\newlabel{lemma-derived-completion-koszul}{{0920}{source}{}{sourcepointer-3}{}}
\newlabel{lemma-derived-completion-koszul@cref}{{[originallemma][0][]0920}{[1][0][]source}{}{}{}}
\makeatother
\title{Derived Nakayama: a derived complete object whose reduction modulo a finitely generated ideal vanishes is zero}
\author{The Stacks Project Authors}
\date{stacks-project:more-algebra.tex:local-snapshot-20261012; GNU FDL 1.2 or later}
\begin{document}
\maketitle
\setcounter{subsection}{1}
\noindent\textit{Editorial note.} All delivered mathematical text is the author's own, copied from the pinned Stacks Project chapter file more-algebra.tex; no statement or proof step was rewritten, reordered or removed. The counted statement is the labeled lemma environment at source lines 26471-26487 (canonical tag 0G1U) and the counted proof is the proof environment at source lines 26489-26501, which immediately follows it and is its own proof.\par\medskip
\noindent\textit{Editorial note.} The counted statement contains the author's slogan recording the name of the result and the author's reference environment recording the history of the statement, including a citation and a URL. The retained preamble defines both environments as comment-type, so their bodies are discarded when the document is typeset and the citation is never evaluated; this is why the card needs no bibliography, and both environments are kept verbatim inside the counted span.\par\medskip
\noindent\textit{Editorial note.} The three references are all to earlier results of this chapter; the kinds recorded are the type words appearing in the labels themselves. All three are navigation stubs with their roles recorded in the reasons.\par\medskip
\noindent\textit{Editorial note.} The preamble span runs to line 90, the current end of the retained preamble, which is exactly what makes the slogan and reference environments comments; the excerpt serves only as preamble, never as a span coordinate.\par\medskip
\noindent\textit{Editorial note.} The source URL is the source-level Stacks Project URL, as required when more than one card is produced from one source; the per-claim tag is carried by claim\_key and by this note.\par\medskip
\noindent\textit{Allowed standard prerequisites (editorial).}\begin{itemize}
\item The Koszul complex on the n-th powers of a generating set is bounded and its cohomology modules are annihilated by a power of the ideal. used at proof lines 26491-26496. Shows that tensoring the object with those complexes vanishes, by the vanishing statement for modules annihilated by a power of the ideal.
\item A derived complete object is the derived limit of its tensor products with the Koszul complexes of the system. invoked at proof lines 26496-26500. Concludes that the object itself is zero once all those tensor products vanish.
\item The special case of the boundedness statement for a derived complete object also gives the result, as the author notes at the start. used at proof line 26490. Provides an alternative shorter proof, after which the direct proof is given.
\end{itemize}\medskip
\ifdefined\hypertarget\hypertarget{sourcepointer-1}{}\fi
\noindent\textit{Source reference.} Original lemma 0H83: sources/stacks-more-algebra/pinned-original-source.tex, source lines 26437--26437. The reference is mentioned at proof line 26490 as an alternative shorter proof; the direct argument actually given is inside the counted span. The lemma itself is the subject of a separate card of this lane with disjoint coordinates.\par\medskip
\ifdefined\hypertarget\hypertarget{sourcepointer-2}{}\fi
\noindent\textit{Source reference.} Original lemma 0G1T: sources/stacks-more-algebra/pinned-original-source.tex, source lines 24155--24155. The reference is applied at proof lines 26494-26496 to kill the tensor products with the Koszul complexes; the conclusion drawn from it is inside the counted span.\par\medskip
\ifdefined\hypertarget\hypertarget{sourcepointer-3}{}\fi
\noindent\textit{Source reference.} Original lemma 0920: sources/stacks-more-algebra/pinned-original-source.tex, source lines 26388--26388. The reference describes derived completeness through the derived limit of the Koszul tensor products, used at proof lines 26496-26498; the final vanishing conclusion is inside the counted span.\par\medskip
\begin{lemma}
\label{lemma-derived-complete-zero-bis}
\begin{slogan}
Derived Nakayama
\end{slogan}
\begin{reference}
A related result is \cite[Proposition 6.5]{Dwyer-Greenlees}.
The derived Nakayama lemma can for example be found in Bhatt's 3rd lecture
on Prismatic cohomology at Columbia University in Fall 2018
as Section 2 property (2). Leonid Positselski proposed a proof in
\url{https://mathoverflow.net/a/331501}. However, we follow the proof
suggested by Anonymous in the comments.
\end{reference}
Let $I$ be a finitely generated ideal of a ring $A$.
Let $K$ be a derived complete object of $D(A)$.
If $K \otimes_A^\mathbf{L} A/I = 0$, then $K = 0$.
\end{lemma}

\begin{proof}
This is a special case of Lemma \ref{lemma-derived-complete-bounded}
but we'll also give a direct proof.
Choose generators $f_1, \ldots, f_r$ of $I$. Denote $K_n$ the
Koszul complex on $f_1^n, \ldots, f_r^n$ over $A$. Recall that
$K_n$ is bounded and that the cohomology modules of $K_n$ are
annihilated by $f_1^n, \ldots, f_r^n$ and hence by $I^{nr}$.
By Lemma \ref{lemma-derived-vanishing-mod-I}
we see that $K \otimes_A^\mathbf{L} K_n = 0$.
Since $K$ is derived complete by Lemma \ref{lemma-derived-completion-koszul}
we have $K = R\lim K \otimes_A^\mathbf{L} K_n = 0$
as desired.
\end{proof}

\end{document}
